\documentclass[11pt]{article}
\usepackage[T1]{fontenc}\usepackage{lmodern}
\usepackage[margin=1in]{geometry}
\usepackage{amsmath,amssymb,amsthm,microtype,booktabs}
\usepackage[hidelinks]{hyperref}
\newtheorem{theorem}{Theorem}[section]\newtheorem{lemma}[theorem]{Lemma}
\newtheorem{proposition}[theorem]{Proposition}
\newcommand{\R}{\mathbb R}\newcommand{\Z}{\mathbb Z}
\newcommand{\ran}{\operatorname{ran}}\newcommand{\rank}{\operatorname{rank}}
\newcommand{\lean}[2]{\par\smallskip\begingroup\small\noindent\textit{Source correspondence:}
\texttt{\detokenize{#1}}\par\noindent\texttt{\detokenize{#2}}\par\endgroup\smallskip}
\title{Gauge pinning and a rational boundary Schur gap\\on a one-sided four-dimensional cell}
\author{Project Stardust}\date{October 29, 2026}
\begin{document}\maketitle
\begin{abstract}
We reconstruct a finite plaquette-incidence problem with 48 vertices,
112 edges and 72 faces. The 32 interior edges have a one-dimensional gauge
kernel; pinning one edge gives an invertible 31-coordinate interior Gram
matrix without changing its curvature range. Eliminating those interior
variables produces a boundary Schur form whose kernel is exactly the
46-dimensional boundary-gradient space. On its 34-dimensional orthogonal
complement we prove a computer-assisted rational lower bound $0.69639$ and
give an exact nonzero Rayleigh witness strictly below $0.6964$. Complete
geometric and linear-algebra arguments identify the physical norm, cycle
compression and residual certificate. The finite arithmetic is supplied as
explicit integer tables with an independent standard-library checker.
A periodic gauge counterexample and a perturbation cross-term test delimit
which invertibility and stability conclusions transfer. This is a finite
quadratic construction, not a nonlinear gauge, multiscale or continuum
mass-gap theorem; the general Schur and discrete exactness mechanisms are
classical.
\end{abstract}

\section{The actual finite problem}
Let $e_0,e_1,e_2,e_3$ be the coordinate unit vectors in $\R^4$. Coordinates
and directions in this paper are zero based, matching the formal source.
The base vertices are $V_0=\{0,1\}^4$. The interior oriented edges are
\[
 E_I=\{(v,i):v\in V_0,\ v_i=0\},
\]
with source $v$ and target $v+e_i$. There are $4\cdot2^3=32$ such edges.
For each $v\in V_0$ with $v_j=1$, add a cap vertex $v+e_j$ and an outward
edge from $v$ to $v+e_j$. These cap vertices are distinct because they have
exactly one coordinate equal to two. There are 32 caps and outward edges.

For every interior edge $(v,i)$ and every $j\ne i$ with $v_j=1$, add a
tangential cap edge from $v+e_j$ to $v+e_i+e_j$. These give 48 tangential
edges: choose the ordered distinct directions $(i,j)$ in 12 ways and the
remaining two binary coordinates in four ways. Let $E_B$ consist of the
32 outward and 48 tangential edges. There are 80 boundary edges, 48 total
vertices $V$, and 112 total edges $E=E_I\sqcup E_B$.

There are two sorts of faces. The 24 inner squares have base $v\in V_0$,
directions $i<j$ and $v_i=v_j=0$. The 48 ribbons correspond to the triples
$(v,i,j)$ used for tangential edges. A ribbon traverses the interior edge
$(v,i)$, the outward edge at $v+e_i$ in direction $j$, the negative
tangential cap edge and the negative outward edge at $v$ in direction $j$.
Thus each is a genuine oriented square. Put $F=F_I\sqcup F_R$, of size 72.

\begin{center}
\begin{tabular}{lrr}
\toprule
Object & Construction & Count\\\midrule
Base/cap vertices & $16+32$ &48\\
Interior edges & $4\cdot8$ &32\\
Outward/tangential boundary edges & $32+48$ &80\\
Inner squares/ribbons & $24+48$ &72\\\bottomrule
\end{tabular}
\end{center}
\lean{OneSidedBlockIncidence}{edges_card; innerFaces_card; upperFaces_card}
\lean{OneSidedCapComplex}{boundary_card; fullVertex_card; fullEdge_card; fullFace_card}

All edge and face spaces have their ordinary Euclidean norm, with every
edge and face counted once. For a potential $\phi\in\R^V$, define
$d_0\phi(e)=\phi(\operatorname{target}e)-\phi(\operatorname{source}e)$.
For $f\in\R^E$, let $d_1f\in\R^F$ be the signed sum around each oriented
face. Telescoping around a square gives $d_1d_0=0$.
The energy is exactly $\|d_1f\|^2$, with no weighting or gauge penalty
implicitly added.

The source also identifies this cap description with the following direct
plaquette rule in $\{0,1,2\}^4$: retain all unit squares incident to at least
one edge in $E_I$, and retain their edges. It obtains the same 72 faces and
112 edges, up to row orientation. Reversing a row changes its sign and not
the energy. The accompanying checker independently reconstructs this rule
and verifies every incidence entry in the certificate.
\lean{OneSidedCapComplex}{embedFace_bijective; d1_d0}
\lean{OneSidedSourceRows}{row_transport; energy_transport}

\section{Exactness by paths, not an assumed gauge quotient}
\begin{lemma}[Every closed field is a gradient]\label{lem:exact}
For this cell, $d_1f=0$ if and only if $f=d_0\phi$ for some potential
$\phi$. The kernel of $d_0$ consists exactly of constant potentials.
Consequently $\rank d_0=47$ and $\rank d_1=65$.
\end{lemma}
\begin{proof}
For a closed field, first restrict it to the base cube. Define $\phi(v)$
by summing its edge values from the origin to $v$, flipping the nonzero
coordinates in increasing order. Any two orders of these coordinate flips
can be related by adjacent interchanges. Interchanging flips in distinct
directions changes the path sum by the signed curl around their inner
square, which is zero. Thus the sum is independent of their order.

For any base edge $(v,i)$ with $v_i=0$, a path to $v+e_i$ may therefore be
chosen as a path to $v$ followed by that edge. It follows that
$\phi(v+e_i)-\phi(v)=f(v,i)$. This proves the gradient identity on every
interior edge. For a cap vertex $v+e_j$, define
$\phi(v+e_j)=\phi(v)+f(v,j)$, where this last edge is the outward one.
This makes the gradient agree with $f$ on all outward edges.
The zero-curl equation on each ribbon then forces agreement on its
tangential edge as well. These are all edges, so $f=d_0\phi$.
The converse is the already established square telescoping identity.

The base cube is connected by coordinate paths, and each cap vertex is
joined to its base vertex by an outward edge. Hence the full graph is
connected. If $d_0\phi=0$, potential differences vanish along every edge,
so all potential values are the same. Constants conversely have zero
gradient. Rank--nullity gives $\rank d_0=48-1=47$.
Exactness makes $\dim\ker d_1=47$, so $\rank d_1=112-47=65$.
\end{proof}
\lean{CubePathPoincare}{path_stokes; cube_poincare}
\lean{OneSidedCapPoincare}{d0_recover; closed_iff_gradient}
\lean{OneSidedGradientRanks}{gradient_kernel_constant; gradient_rank; curl_rank}
The formal proof implements an ordered-path potential and its exact Stokes
error formula. The adjacent-swap proof above reconstructs that same finite
path mechanism. No cohomological exactness is assumed merely from the shape
of a diagram.

\section{The interior gauge and a legitimate pin}
Write $A:\R^{E_I}\to\R^F$ for the restriction of $d_1$ to fields that
vanish on the boundary, and $B:\R^{E_B}\to\R^F$ for its boundary restriction.
Thus the full curvature is $Af+B\eta$. The interior energy is
\begin{equation}\label{eq:intenergy}
 \|Af\|^2=\sum_{\text{inner squares}}(\operatorname{curl}f)^2
                  +\sum_{(v,i,j)\text{ ribbons}}f(v,i)^2.
\end{equation}
The second sum occurs because only the interior edge of a ribbon can
contribute when the boundary field is zero.

Let $g\in\R^{E_I}$ equal one on the four edges leaving the origin and
zero on the other interior edges. Its squared norm is four. It is a full
gradient: take the potential $\psi$ equal to $-1$ at the base origin and
zero at every other vertex. None of the boundary edges is incident to
the origin, so $d_0\psi$ equals $g$ on the interior and zero on the boundary.
In particular $Ag=0$.

\begin{lemma}[Exact interior kernel]\label{lem:intkernel}
The kernel of $A$ is precisely $\R g$; hence $\rank A=31$.
\end{lemma}
\begin{proof}
If $Af=0$, every square and ribbon summand in \eqref{eq:intenergy} is zero.
An interior edge $(v,i)$ away from the origin has $v_i=0$ and some other
coordinate $v_j=1$. It belongs to a ribbon, which forces $f(v,i)=0$.
Only the four origin edges can remain. The inner square at the origin
in directions $i,j$ now has curl $f(0,i)-f(0,j)$ because its other two
edges are away from the origin. Thus all four remaining values agree,
so $f=cg$. Conversely $Ag=0$, proving the kernel statement.
Rank--nullity gives $32-1=31$.
\end{proof}
\lean{OneSidedBlockKernel}{kernel_classification; energy_zero_iff}
\lean{OneSidedBlockRank}{kernel_span; incidence_rank; gauge_is_gradient}

Choose the pole edge $e_*=(0,0)$. Pin its value to zero and let
$L:\R^{31}\to\R^{72}$ be the resulting interior incidence matrix.
For any unpinned $f$, subtract $f(e_*)g$. This does not change $Af$ and
makes the pole value zero. Thus
\begin{equation}\label{eq:range}
 \ran L=\ran A.
\end{equation}
If a pinned field has zero curvature, Lemma~\ref{lem:intkernel} says it is
$cg$, and its zero pole value forces $c=0$. Therefore $L$ is injective and
$C=L^TL$ is positive definite. It has an actual inverse: in finite dimension
a positive-definite symmetric matrix is injective and hence bijective.
\lean{OneSidedPinnedInterior}{lift_pin; pinned_range; pinned_injective; localGram_isUnit}
The unpinned matrix $A^TA$ is singular. Only the proved pin or a separately
specified generalized inverse can justify elimination. Equation~\eqref{eq:range}
will allow us to check normal equations in the unpinned coordinates without
ever inverting that singular matrix.

\section{The boundary gradient and its physical complement}
Let $G:\R^{48}\to\R^{80}$ be the restriction of $d_0$ to boundary edges.
Its range is the boundary-gradient space, and define the physical boundary
space by
\[
 \mathcal P=(\ran G)^\perp=\ker G^T.
\]
This is a finite real linear definition. It is not an identification with
a nonlinear gauge-invariant quantum state space.

\begin{lemma}[The complete boundary gauge]\label{lem:bgauge}
The kernel of $G$ is the two-dimensional space of potentials
$\phi=a\mathbf1+c\psi$, where $\psi$ is the origin potential defined
above. Hence $\rank G=46$ and $\dim\mathcal P=34$.
\end{lemma}
\begin{proof}
Both constants and $\psi$ have zero boundary gradient. They are linearly
independent, since $\psi$ is zero at cap vertices and nonzero at the origin.
Conversely suppose $G\phi=0$. Then $d_0\phi$ has only interior components,
and its full curl vanishes. Its interior restriction $f$ satisfies $Af=0$,
so $f=cg$ by Lemma~\ref{lem:intkernel}. As full edge fields,
$d_0\phi=c\,d_0\psi$. Thus $d_0(\phi-c\psi)=0$.
Lemma~\ref{lem:exact} makes $\phi-c\psi$ constant, proving the complete
kernel classification. Rank--nullity gives $48-2=46$ for $G$, and ordinary
Euclidean orthogonal decomposition gives $\dim\mathcal P=80-46=34$.
\end{proof}
\lean{OneSidedBoundaryGauge}{boundary_kernel; boundary_gradient_rank; boundary_quotient_finrank}
\lean{OneSidedBoundaryCycles}{physical_finrank; divergence_zero_iff}
Boundary gauge removal therefore costs 46 dimensions, not the single
interior gauge dimension. Subtracting one dimension from both spaces would
describe a different boundary problem.

\section{Schur elimination and its exact kernel}
For a boundary field $\eta$, define the pinned response, residual and Schur
matrix by
\begin{align*}
 y_\eta&=-(L^TL)^{-1}L^TB\eta,\\
 r_\eta&=Ly_\eta+B\eta,\\
 S&=B^TB-B^TL(L^TL)^{-1}L^TB.
\end{align*}
The normal equation is $L^Tr_\eta=0$. In particular $r_\eta$ is orthogonal
to every pinned interior curvature.

\begin{theorem}[Completed square and physical kernel]\label{thm:schur}
For all $y\in\R^{31}$ and $\eta\in\R^{80}$,
\begin{equation}\label{eq:square}
 \|Ly+B\eta\|^2=\eta^TS\eta+\|L(y-y_\eta)\|^2,
 \qquad \eta^TS\eta=\|r_\eta\|^2.
\end{equation}
The minimizing pinned $y$ is uniquely $y_\eta$, $S$ is symmetric positive
semidefinite, and $\ker S=\ran G$. Thus $\rank S=34$, and its quadratic
form is strictly positive on $\mathcal P\setminus\{0\}$.
\end{theorem}
\begin{proof}
Write $Ly+B\eta=r_\eta+L(y-y_\eta)$. The cross term is zero by the normal
equation. Also $B^Tr_\eta=S\eta$, and
\[
 \eta^TS\eta=\langle B\eta,r_\eta\rangle
 =\langle r_\eta-Ly_\eta,r_\eta\rangle=\|r_\eta\|^2.
\]
This proves \eqref{eq:square}. Injectivity of $L$ gives uniqueness of its
minimizer. The formula for $S$ is symmetric because $L^TL$ and its inverse
are symmetric, and the squared-residual identity proves nonnegativity.

If $\eta\in\ran G$, choose $\phi$ with $G\phi=\eta$. The full gradient
$d_0\phi$ has zero curvature, so some interior completion has
$Af+B\eta=0$. Replace $f$ by its pinned representative using
\eqref{eq:range}. The minimum energy is zero, hence $r_\eta=0$ and
$S\eta=B^Tr_\eta=0$.
Conversely, if $S\eta=0$, then \eqref{eq:square} gives $r_\eta=0$.
The full field comprising $y_\eta$ and $\eta$ is curl free and therefore a
gradient by Lemma~\ref{lem:exact}. Its boundary restriction is $\eta$, so
$\eta\in\ran G$. This proves the kernel equality and the rank.
If a physical $\eta$ has zero energy, its residual is zero and it lies
in $\ran G$ as well as its orthogonal complement; testing it against
itself gives $\eta=0$. This proves strict physical positivity.
\end{proof}
\lean{OneSidedSchurForm}{completed_square; unique_minimizer; schur_posSemidef}
\lean{OneSidedSchurForm}{kernel_schur; schur_rank; physical_positive}
These steps are classical Schur-complement reasoning, applied after a
specific gauge pin has been proved legitimate. Gallier's primary exposition
\cite{gallier}, Proposition~2.1 and Theorem~4.3, supplies the general
positive-definite and singular-block context. Here the kernel and range
conditions have been established for the actual incidence problem.

\section{Physical cycles and the norm that compression retains}
The certificate supplies an integer matrix $C_b\in\Z^{80\times34}$.
Its columns are actual boundary cycles. Two finite identities are crucial:
\begin{equation}\label{eq:cycles}
 G^TC_b=0,\qquad (C_b)_{c_i,j}=\mathbf1_{i=j}
 \quad(0\leq i,j<34),
\end{equation}
where $c_0,\ldots,c_{33}$ are distinct chord-row indices, also supplied.
The first identity puts every column in $\mathcal P$. The identity minor
makes the columns linearly independent: if $C_bz=0$, reading its chord
rows gives $z=0$. Since $\dim\mathcal P=34$, these columns span all of
$\mathcal P$. Every physical field is uniquely $\eta=C_bz$.
\lean{OneSidedBoundaryCycleData}{cycle_balance; chord_identity}
\lean{OneSidedBoundaryCycles}{cycle_injective; cycle_range; physical_representation}

Put
\[
 H=C_b^TC_b,\qquad K=C_b^TSC_b.
\]
Then
\begin{equation}\label{eq:metric}
 \|\eta\|^2=z^THz,\qquad \eta^TS\eta=z^TKz.
\end{equation}
In particular $H$ is positive definite. The cycle coordinates are not
orthonormal; replacing $H$ by the identity would change the Rayleigh
quotient and the claimed physical gap. We keep the generalized quotient
$z^TKz/(z^THz)$ throughout.
\lean{OneSidedSchurCompression}{metric_quadratic; compressed_quadratic; gap_transfer}
The orthogonal split into gradient and divergence-free edge fields is
classical graph Hodge theory. Jiang--Lim--Yao--Ye \cite{hodge} describe this
mechanism in their primary treatment of edge flows, Theorems 4.7--4.8. Their application is
different; no statistical ranking conclusion or universal graph gap is
imported into the present plaquette certificate.

\section{A normal-equation certificate without a singular inverse}
The companion data file contains integer matrices
\[
 X\in\Z^{32\times34},\quad R\in\Z^{72\times34},\quad
 Q\in\Z^{34\times34},\quad H\in\Z^{34\times34},\quad D=193028.
\]
Here $A,B,C_b$ use the explicit source plaquette ordering in that file.
The checked identities are
\begin{align}
 AX+D BC_b&=R,& A^TR&=0,\label{eq:normalcert}\\
 C_b^TC_b&=H,& R^TR&=DQ.\label{eq:gramcert}
\end{align}
All are integer equalities. Their exact arithmetic verification is described
in Section~\ref{sec:checks}. We first prove their mathematical implication.

\begin{lemma}[Identifying the certificate with the original Schur form]
Identities \eqref{eq:normalcert}--\eqref{eq:gramcert} imply $K=Q/D$.
\end{lemma}
\begin{proof}
For any $z\in\R^{34}$, let $\eta=C_bz$ and $f=Xz/D$.
The first equation says $Af+B\eta=Rz/D$. The second says this residual is
orthogonal to $\ran A=\ran L$. The original minimizing residual $r_\eta$
also lies in $B\eta+\ran A$ and is orthogonal to that range.
The difference of these two residuals lies both in $\ran A$ and its
orthogonal complement, so is zero. Therefore $r_\eta=Rz/D$.
Taking its squared norm and using \eqref{eq:gramcert} gives
\[
 z^TKz=\|r_\eta\|^2
 =\frac{z^TR^TRz}{D^2}=\frac{z^TQz}{D}.
\]
Both matrices are symmetric; equality of their quadratic forms for all $z$
implies equality of matrices, by testing coordinate vectors and their
pairwise sums. Hence $K=Q/D$.
\end{proof}
\lean{OneSidedIntegerResponseData}{response_equation; normal_equation}
\lean{OneSidedCertifiedResidual}{normal; compressed_eq_residual_gram; compressed_entry}
\lean{OneSidedBoundaryGap}{compressed_as_data; metric_as_data}
Notice that $X$ uses all 32 unpinned coordinates. Its normal equations are
valid without an inverse of $A^TA$. Uniqueness of the \emph{residual}, not
uniqueness of the unpinned response, is what identifies the certificate.
Row orientation signs in the native cap model cancel in the residual Gram
and normal equations; the checker independently verifies the canonical
lattice row geometry.

\section{The rational lower certificate and a sub-rounded witness}
Define the integer matrix
\begin{equation}\label{eq:Tdata}
 T=25000Q-69639\cdot48257H.
\end{equation}
The certificate supplies a positive integer $s$, positive integers
$w_0,\ldots,w_{33}$ and an integer matrix $V\in\Z^{34\times34}$ such that
\begin{equation}\label{eq:positive}
 sT=V\operatorname{diag}(w_0,\ldots,w_{33})V^T.
\end{equation}
The exact value of $s$ has 2,386 decimal digits. All of it, every weight and
every matrix entry are provided in the accompanying data; no floating-point
factorization or approximate positivity check is used.

\begin{theorem}[Certified physical boundary gap]\label{thm:gap}
For every $\eta\in\mathcal P$,
\[
 \frac{69639}{100000}\|\eta\|^2\leq\eta^TS\eta.
\]
If $\lambda_*$ is the minimum physical Rayleigh quotient, then
\begin{equation}\label{eq:bracket}
 \frac{69639}{100000}\leq\lambda_*
 \leq\frac{729067245}{1046910149}<\frac{1741}{2500}=0.6964.
\end{equation}
\end{theorem}
\begin{proof}
For any $z$, \eqref{eq:positive} gives
$sz^TTz=\sum_k w_k((V^Tz)_k)^2\geq0$. Division by $s>0$ gives
$z^TTz\geq0$. Since $D=4\cdot48257$, direct substitution into
\eqref{eq:Tdata} yields
\[
 K-\frac{69639}{100000}H=\frac{T}{25000D}
                            =\frac{T}{4825700000}.
\]
The denominator is positive. Equations~\eqref{eq:metric} therefore prove
the gap for every physical field.

For the upper comparison, the certificate's explicit vector is
\begin{align*}
 z_*={}&(2000,0,2000,760,2000,2000,-704,-704,9,-1935,\\
 &-704,-704,-1935,9,-1935,9,-130,-130,-1935,-704,\\
 &-704,-1935,-1935,0,0,-56,-56,-56,-56,-130,\\
 &-56,-56,-130,-130)^T.
\end{align*}
The exact integer dot products are
\[
 z_*^THz_*=117410484,\qquad z_*^TQz_*=15782847719760.
\]
Thus $\eta_*=C_bz_*$ is physical and nonzero, and its quotient is
\[
 \frac{\eta_*^TS\eta_*}{\|\eta_*\|^2}
 =\frac{15782847719760}{193028\cdot117410484}
 =\frac{729067245}{1046910149}<\frac{1741}{2500}.
\]
The final comparison is exact cross multiplication of positive integers.
The minimum exists: the unit sphere of the finite-dimensional physical
space is compact and the quadratic form is continuous. Its lower and upper
bounds follow from the proved inequality and this test field.
\end{proof}
\lean{OneSidedBoundaryPositiveData}{scale_positive; weights_positive; weighted_gram_identity}
\lean{OneSidedBoundaryGap}{corrected_eq; compressed_bound; boundary_gap}
\lean{OneSidedBoundaryGapWitness}{integer_norm; integer_gram; rayleigh_value; below_rounded}
The witness quotient is approximately $0.69639906127$; this decimal is only
an orientation aid. In particular the rounded decimal $0.6964$ is
\emph{not} a valid lower bound, even though the certified interval is less
than $10^{-5}$ wide. The rational lower and upper tests have different roles.
No assertion that the witness is an exact eigenvector or minimizer is made.
The compactness statement introducing $\lambda_*$ is a classical human
rephrasing of the mapped quadratic-form bounds, not a new formal declaration.

\section{What the finite arithmetic checks}\label{sec:checks}
This theorem is computer assisted at its finite integer-certificate stage.
The geometric exactness, gauge pin, completed square, norm identification
and deduction from a positive Gram certificate were proved above. The
remaining finite input is fully specified by
\texttt{certificate-data.json} and checked by \texttt{verify\_certificate.py}.
The JSON contains ordinary integer arrays and their original source hashes.
Its SHA-256 digest, using UTF-8 text with LF newlines, is
\begin{quote}\small\texttt{b1a7d776078753d95a18c451f4536b6cc4aeea2b12142fedaf8e66b972d01078}.\end{quote}
It includes the interior and boundary edge coordinates, source face
coordinates, boundary endpoint indices, chord indices, $A,B,C_b,X,R,H,Q,T$,
$s,w,V$ and $z_*$. Directions and array indices are zero based.

The checker uses only Python's standard library and unbounded integer
arithmetic, together with exact rational arithmetic for the last quotient.
Its logical operations are the following finite computations.
\begin{enumerate}
\item Generate all unit squares in $\{0,1,2\}^4$ incident to one of the
32 interior edges; verify the 72-face, 112-edge, 48-vertex geometry and the
32/80 split against the listed coordinate arrays.
\item For every face, assign the four coefficients $+1,+1,-1,-1$ along
its canonical boundary. Verify every entry of $A,B$ against those rows.
\item Construct the boundary gradient matrix from its actual source and
target vertices. Verify both equations in \eqref{eq:cycles}, including all
34 chord rows.
\item Multiply the integer arrays to verify every entry of
\eqref{eq:normalcert} and \eqref{eq:gramcert}.
\item Verify \eqref{eq:Tdata}, check $s>0$ and every $w_k>0$, and verify
all 1,156 entries of \eqref{eq:positive} by ordinary integer multiplication
and addition.
\item Verify the two witness dot products and the rational comparisons
in \eqref{eq:bracket} by exact integer or fraction arithmetic.
\end{enumerate}
Each matrix product uses the defining formula
$(MN)_{ij}=\sum_kM_{ik}N_{kj}$. Thus no eigenvalue solver, numerical tolerance,
symbolic positivity oracle or unproved matrix inversion is hidden in the
finite certificate. These checks pass independently of the source's
kernel-checked integer identities. The companion receipt records the data
hash and the actual successful check; it is not a substitute for the
analytic and geometric implications of those identities.

\section{A local pin is not global unpinned invertibility}
There is a concrete periodic counterexample to transferring a local pinned
statement to simultaneous elimination of all unpinned block interiors.
Let the periodic vertices be $(\Z/4\Z)^4$, with positive coordinate edges.
Partition them into side-two blocks by the coordinate labels
$\lfloor v_i/2\rfloor$. Call an edge interior if its two endpoints have
the same four block labels.
Define a periodic potential
\[
 \Phi(v)=\begin{cases}1,&v_0\in\{1,2\},\\0,&v_0\in\{0,3\}.
 \end{cases}
\]
Let $f_i(v)=\Phi(v+e_i)-\Phi(v)$. It vanishes in directions $i\ne0$.
In direction zero, the values for $v_0=0,1,2,3$ are respectively
$1,0,-1,0$. The nonzero edges run from zero to one and from two to three,
so they are interior edges. The crossing edges from one to two and from
three to zero have zero value. Hence $f$ is nonzero and entirely supported
on the simultaneously retained interiors.

Every plaquette curl of $f$ is zero, because it is a gradient and the
coordinate shifts commute. If $A_{\rm per}$ denotes the full periodic
plaquette incidence restricted to all these interior edges, its nonzero
interior field $f_I$ satisfies $A_{\rm per}f_I=0$. Consequently
$A_{\rm per}^TA_{\rm per}f_I=0$. It has no left inverse and no strictly
positive unpinned coercivity constant: either assertion applied to $f_I$
would force it to vanish or give positive norm bounded by zero.
\lean{PeriodicInteriorGaugeKernel}{gauge_boundary_zero; witness_nonzero; incidence_witness_zero}
\lean{PeriodicInteriorGaugeKernel}{hessian_witness_zero; hessian_no_left_inverse}
This is compatible with the local result. The one-sided pinned problem
and the unrestricted periodic all-block interior matrix are different
operators. The counterexample does not exclude a carefully defined global
gauge quotient or a different pinned global estimate; those would require
their own domain, kernel and coercivity proofs.

\section{A separate perturbation boundary: missing linear loss}
Even finite positive blocks do not allow a general perturbation estimate
that drops coupling cross terms. For a scalar block
\[
 H=\begin{pmatrix}a&b\\b&c\end{pmatrix},\qquad c>0,
 \qquad \mathcal S(a,b,c)=a-b^2/c,
\]
changing $b$ to $b+d$ at $c=1$ gives the exact identity
\begin{equation}\label{eq:cross}
 \mathcal S(a,b+d,1)=\mathcal S(a,b,1)-2bd-d^2.
\end{equation}
At $a=2,b=1,d=1/10$, both blocks are positive definite:
their Schur values are $1$ and $79/100$. But
\[
 \frac{79}{100}<1-4(1/10)^2=\frac{96}{100}.
\]
Thus a proposed loss of only $4d^2$ relative to an arbitrary nonzero
reference coupling is already false in dimension two. More generally
$2d+d^2>4d^2$ for $0<d<2/3$, so arbitrarily small positive perturbations
have the same defect. For $|d|\leq\varepsilon$, \eqref{eq:cross} instead
gives the valid bound
\[
 \mathcal S(a,b+d,1)\geq\mathcal S(a,b,1)
                              -2|b|\varepsilon-\varepsilon^2.
\]
\lean{ManuscriptSchurDiagnostics}{fixed_interior_increment; general_stability_counterexample}
\lean{ManuscriptSchurDiagnostics}{small_coupling_counterexample; fixed_interior_lower_bound}

For comparison, quadratic loss is justified at a zero reference coupling.
If a real Hilbert-space block energy is
\[
 \mathcal E(x,y)=ax^2+2x\langle b,y\rangle+\langle Cy,y\rangle,
\]
with $a\geq a_0-\delta$, $\|b\|\leq\varepsilon$ and $C\geq mI$, $m>0$,
then completing the scalar square yields
\begin{align*}
 \mathcal E(x,y)&\geq(a_0-\delta)x^2
                    -2|x|\varepsilon\|y\|+m\|y\|^2\\
 &\geq\left(a_0-\delta-\frac{\varepsilon^2}{m}\right)x^2.
\end{align*}
If instead $C_0\geq m_0I$ and $\|C-C_0\|\leq\eta\leq m_0/2$,
Cauchy--Schwarz gives
$\langle Cy,y\rangle\geq(m_0-\eta)\|y\|^2\geq(m_0/2)\|y\|^2$.
The same proof gives loss $\delta+2\varepsilon^2/m_0$.
These are conditional perturbation statements, with a zero-coupling
reference and proved interior coercivity. They do not automatically
transport the numerical cell gap to nonlinear fields or other geometries.
\lean{CoerciveSchurStability}{operator_quadratic_perturbation; operator_robust_lower}

\section{Provenance, significance and remaining uniformity questions}
The mathematics is reconstructed from frozen verified commit
\begin{quote}\small\texttt{bab688bb6140cd6316c80680095a914ea0d346db}.\end{quote}
The exact declaration map and original source hashes accompany the paper.
The independent checker verifies a literal extraction of the frozen integer
tables against geometry; no mathematical research source, configuration or
original manuscript was altered. The source's complete build and axiom audit
receipt was checked by its evidence hashes; no new research build is claimed.

Finite exactness, gauge gradients, Schur complements and positive weighted
Gram certificates use classical mechanisms. The primary comparisons above
establish that context, not novelty of general methods. Priority for the
specific one-sided geometry and exact narrow bracket remains unresolved.
The result is a reproducible geometry-specific finite quadratic estimate,
including an explicit correction to rounding and a sharp distinction between
local pinning and a global unpinned matrix.

There is no varying block size, lattice spacing, volume or continuum limit
in the theorem. It neither constructs a four-dimensional quantum field nor
identifies a physical correlation or Hamiltonian spectrum. The linear
boundary complement is not silently substituted for a nonlinear gauge
quotient. Extending the estimate would require exact compatible global gauge
conditions, control of interactions and approximation errors, and estimates
uniform in the particular scale and volume parameters. The proved finite
gap discharges none of these separate analytic obligations.

\begin{thebibliography}{9}
\bibitem{gallier} J. Gallier, \emph{The Schur complement and symmetric positive
semidefinite (and definite) matrices}, author notes, August 24, 2019.
\href{https://www.cis.upenn.edu/~jean/schur-comp.pdf}{Primary author PDF}.
\bibitem{hodge} X. Jiang, L.-H. Lim, Y. Yao and Y. Ye,
\emph{Statistical ranking and combinatorial Hodge theory}, primary preprint,
\href{https://arxiv.org/abs/0811.1067}{arXiv:0811.1067}.
\end{thebibliography}
\end{document}
